% Generated draft: requires mathematical and source review.
\documentclass[11pt]{article}
\usepackage[margin=25mm]{geometry}
\usepackage{fontspec}
\setmainfont{Latin Modern Roman}
\usepackage{amsmath,amssymb,mathtools,unicode-math}
\setmathfont{Latin Modern Math}
\usepackage{enumitem,needspace,xurl,hyperref}
\setlength{\parindent}{0pt}
\setlength{\parskip}{5pt}
\setlength{\emergencystretch}{3em}
\newcommand{\N}{\mathbb{N}}
\newcommand{\Z}{\mathbb{Z}}
\newcommand{\Q}{\mathbb{Q}}
\newcommand{\R}{\mathbb{R}}
\newcommand{\C}{\mathbb{C}}
\newcommand{\F}{\mathbb{F}}
\newcommand{\A}{\mathbb{A}}
\newcommand{\PP}{\mathbb{P}}
\newcommand{\E}{\mathbb{E}}
\newcommand{\eps}{\varepsilon}
\newcommand{\dd}{\,\mathrm{d}}
\begin{document}
\section*{UnsolvedMath ID: 20001977}
\textbf{Source problem code:} catalogue entry
\subsection{Definitions and mathematical statement}
R. Goldin's Question 10.3 asks for which symplectic toric orbifolds [LT97] a preceding surjectivity statement holds over \(\mathbb Z\), suggests an orbifold-cohomology analogue, and asks about \(\mathbb Z/p\mathbb Z\). The precise preceding map and meaning of orbifold cohomology are not reproduced. The attached Comment 10.4 is attributed to E. Lerman but truncates at ``Note that in''.

The paired aid treats the standard positive-weight circle presentation. Let \(w_1,\ldots,w_n\) be positive integers and let \(S^1\) act on \(\mathbb C^n\) by
\[
 e^{it}(z_1,\ldots,z_n)=(e^{iw_1t}z_1,\ldots,e^{iw_nt}z_n),
 \qquad W=\prod_{j=1}^n w_j.
\]
A positive moment level is equivariantly diffeomorphic to \(S^{2n-1}\); the quotient is the weighted projective orbifold \(\mathbb P(w_1,\ldots,w_n)\), allowing the given presentation even when ineffective. For a coefficient ring \(R\), define
\[
 \kappa_R:H^*_{S^1}(\mathbb C^n;R)=R[u]
       \longrightarrow H^*_{S^1}(S^{2n-1};R),\qquad \deg u=2.
\]
The target is Borel/Haefliger orbifold cohomology of this quotient presentation. It is not automatically the ordinary cohomology of the coarse space, or Chen--Ruan twisted-sector cohomology.

The aid reports integral surjectivity with
\[
 H^*_{S^1}(S^{2n-1};\mathbb Z)
 \cong\mathbb Z[u]/(Wu^n),\qquad
 \ker\kappa_{\mathbb Z}=(Wu^n).
\]
For a prime \(p\), its reported exact criterion is
\[
 \kappa_{\mathbb F_p}\text{ is surjective}
 \quad\Longleftrightarrow\quad p\nmid W.
\]
If \(p\nmid W\), the target is \(\mathbb F_p[u]/(u^n)\). If \(p\mid W\), the equivariant Euler class vanishes modulo \(p\), and the target has a nonzero odd class in degree \(2n-1\) that cannot be in the image of the even-graded source. No exterior-algebra presentation of this odd class is asserted, especially in characteristic two. Tensoring the integral quotient ring alone need not compute the full mod-\(p\) cohomology, because torsion can contribute additional classes.

In particular \(\mathbb P(1,m)\) is an effective family for \(m\geq1\), integrally surjective in this Borel sense but nonsurjective modulo every prime dividing \(m\). This provides a coefficient-sensitive special case, not a classification for all symplectic toric orbifolds or all meanings of orbifold cohomology.
\subsection{Short English statement}
For the weighted-projective Borel/Haefliger model, the supplied aid gives integral Kirwan surjectivity and a sharp mod-prime criterion: surjectivity fails precisely at primes dividing the weight product, because new odd cohomology appears. The original broader toric and alternative-orbifold-cohomology questions remain separately specified tasks.
\subsection{Sources}
\begin{itemize}
\item Local source record: catalogue entry, catalogue entry; canonical \texttt{aim-geometry-notes.json}, \texttt{notes[314]}, source item 10.3. Title: ``Integral and mod-p Kirwan maps for weighted projective orbifolds.''
\item AIM workshop list \emph{Moment maps and surjectivity in various geometries}, Question 10.3 attributed to R. Goldin and truncated Comment 10.4 attributed to E. Lerman; supplied reference \url{https://aimath.org/WWN/momentmaps/momentmaps.pdf}. This draft uses the supplied text without accessing a PDF.
\item Original reference [LT97] for symplectic toric orbifolds; no full entry is supplied. The weighted-projective coefficient criterion is a candidate theorem of research attempt 1, without a separate authoritative citation in the record.
\item Registry dated 14 August 2026: ``partially\_solved,'' ``PARTIAL-PROGRESS,'' statement status ``exact''; attempt partial, rigor ``proved,'' novelty confidence low. The machine-generated, non-peer-reviewed aid requires independent verification before citation.
\end{itemize}
\end{document}
